\section{Class P --- Binary Arith --- Div}
%%% generated by the blueprint pipeline from TCSlib.Complexity.ClassP.BinaryArith.Div
\subsection*{Class P --- Binary Arith --- Div: Introduction}

Declarations of the Lean module \texttt{TCSlib.Complexity.ClassP.BinaryArith.Div}.

\subsection*{Class P --- Binary Arith --- Div: Results}

\begin{theorem}[Polynomial-Time Zero Test]
\lean{Complexity.BinArith.polyTimeComputable_isZero}
\label{Complexity.BinArith.polyTimeComputable_isZero}
\leanok
\uses{Complexity.BinArith.bitsVal_eq_zero_iff, Complexity.BinArith.polyTimeComputable_any, Complexity.LogProg.bitsVal, Complexity.PolyTimeComputable, Complexity.polyTimeComputable_not}
\difficulty{4}
The function sending a finite Boolean word $w$ to $[\mathsf{true}]$ when its value, read
with the first bit as the least significant bit, is zero, and to $[\mathsf{false}]$
otherwise, is polynomial-time computable.
\end{theorem}

\begin{definition}[Binary representation of a quotient]
\lean{Complexity.BinArith.div}
\label{Complexity.BinArith.div}
\leanok
\uses{Complexity.LogProg.bitsVal, Complexity.pairFstD, Complexity.pairSndD}
For a Boolean list $z$, extract its two paired bit lists and interpret them as natural
numbers $a$ and $b$. The definition assigns to $z$ the binary digit list of the
natural-number quotient $a/b$, using the natural-number convention when $b=0$.
\end{definition}

\begin{definition}[Binary representation of a remainder]
\lean{Complexity.BinArith.mod}
\label{Complexity.BinArith.mod}
\leanok
\uses{Complexity.LogProg.bitsVal, Complexity.pairFstD, Complexity.pairSndD}
For a Boolean list $z$, let $a$ and $b$ be the natural numbers decoded from its first
and second paired components, respectively. The definition assigns the binary digit list
of the natural-number remainder of $a$ upon division by $b$, using natural-number
remainder when $b=0$.
\end{definition}

\begin{theorem}[Division of encoded binary numbers]
\lean{Complexity.BinArith.div_pairEncode}
\label{Complexity.BinArith.div_pairEncode}
\leanok
\uses{Complexity.BinArith.div, Complexity.LogProg.bitsVal_bits, Turing.pairEncode}
\difficulty{1}
For all natural numbers $a,b$, dividing the self-delimiting encoding of their binary
expansions produces the binary expansion of $a/b$, where the quotient is integer
division.
\end{theorem}

\begin{theorem}[Remainder of encoded natural numbers]
\lean{Complexity.BinArith.mod_pairEncode}
\label{Complexity.BinArith.mod_pairEncode}
\leanok
\uses{Complexity.BinArith.mod, Complexity.LogProg.bitsVal_bits, Turing.pairEncode}
\difficulty{1}
For natural numbers $a$ and $b$, let their binary expansions be paired by doubling each
bit of the expansion of $a$, appending the separator $01$, and then appending the
expansion of $b$. Applying the operation that decodes the pair, reads the two binary
values, and returns the binary expansion of their remainder yields the binary expansion
of $a \bmod b$.
\end{theorem}

\begin{definition}[Binary division step]
\lean{Complexity.BinArith.divStep}
\label{Complexity.BinArith.divStep}
\leanok
\uses{Complexity.BinArith.le, Complexity.BinArith.sub, Complexity.isNilB, Complexity.pairFstD, Complexity.pairSndD, Turing.pairEncode}
Represent a state $S$ by nested pair encodings, with outer first component $a$, and
outer second component encoding a triple $(d,R,q)$; write $d=[h]\mathbin{+
egthinspace+}t$, where $h$ is its first Boolean digit, taken to be false if $d$ is
empty. If $d$ is empty, leave $S$ unchanged. Otherwise compare $a$ with the binary
number represented by $[h]\mathbin{+
egthinspace+}R$. If $a$ is no greater, return the nested encoding with components $a$,
$t$, the binary subtraction of $[h]\mathbin{+
egthinspace+}R$ by $a$, and the digit true prepended to $q$; if $a$ is greater, return
the encoding with components $a$, $t$, the unchanged $[h]\mathbin{+
egthinspace+}R$, and the digit false prepended to $q$.
\end{definition}

\begin{theorem}[Polynomial-time computability of the division step]
\lean{Complexity.BinArith.polyTimeComputable_divStep}
\label{Complexity.BinArith.polyTimeComputable_divStep}
\leanok
\uses{Complexity.BinArith.divStep, Complexity.BinArith.polyTimeComputable_cons, Complexity.BinArith.polyTimeComputable_le, Complexity.BinArith.polyTimeComputable_slot2, Complexity.BinArith.polyTimeComputable_slot3, Complexity.BinArith.polyTimeComputable_slot4, Complexity.BinArith.polyTimeComputable_sub, Complexity.PolyTimeComputable, Complexity.PolyTimeComputable.comp, Complexity.PolyTimeComputable.pairEncode, Complexity.pairFstD, Complexity.pairSndD, Complexity.polyTimeComputable_const, Complexity.polyTimeComputable_headD, Complexity.polyTimeComputable_id, Complexity.polyTimeComputable_isNil, Complexity.polyTimeComputable_ite, Complexity.polyTimeComputable_pairFstD, Complexity.polyTimeComputable_tail}
\difficulty{5}
The function that maps a bit string to the result of one division-step update, according
to its prescribed cases, is polynomial-time computable.
\end{theorem}

\begin{definition}[Binary division initialization]
\lean{Complexity.BinArith.divInit}
\label{Complexity.BinArith.divInit}
\leanok
\uses{Complexity.pairFstD, Complexity.pairSndD, Turing.pairEncode}
For a list $z$ of Boolean digits, initialize a nested pair encoding whose outer
component is the second decoded component of $z$, whose next component is the reverse of
the first decoded component, and whose final two components are empty lists.
\end{definition}

\begin{definition}[Division invariant]
\lean{Complexity.BinArith.DivInv}
\label{Complexity.BinArith.DivInv}
\leanok
\uses{Complexity.LogProg.bitsVal, Complexity.pairFstD, Complexity.pairSndD, Turing.pairEncode}
For lists $z$ and $S$ of Boolean digits, an invariant holds if there are lists $R$, $r$,
and $q$ such that $S$ is the nested pair encoding of the second decoded component of
$z$, followed by $R$, $r$, and $q$; both $r$ and $q$ have length at most the length of
the first decoded component of $z$ minus the length of $R$. In addition, if the natural
number represented by the second decoded component is positive, then
$(\operatorname{val}(q)\operatorname{val}(b)+\operatorname{val}(r))2^{|R|}+\operatorname{val}(\operatorname{reverse}(R))=\operatorname{val}(a)$
and $\operatorname{val}(r)<\operatorname{val}(b)$, where $a$ and $b$ are the natural
numbers represented by the first and second decoded components of $z$.
\end{definition}

\begin{theorem}[Initial division invariant]
\lean{Complexity.BinArith.divInv_init}
\label{Complexity.BinArith.divInv_init}
\leanok
\uses{Complexity.BinArith.DivInv, Complexity.BinArith.divInit, Complexity.LogProg.bitsVal}
\difficulty{1}
For every bit string $z$, the initial state formed by pairing the second projection of
$z$ with the pairing of the reverse of its first projection and two empty strings
satisfies the following invariant: there exist bit strings $R,r,q$ such that the state
is the nested pairing of the second projection of $z$, $R$, and the pair $(r,q)$;
$|r|+|R|$ and $|q|+|R|$ are each at most the length of the first projection of $z$; and,
if the value of the second projection is positive, then
$(\operatorname{val}(q)\operatorname{val}(\pi_2
z)+\operatorname{val}(r))2^{|R|}+\operatorname{val}(R^{\mathrm{rev}})=\operatorname{val}(\pi_1
z)$ and $\operatorname{val}(r)<\operatorname{val}(\pi_2 z)$. Here bit strings are valued
in little-endian binary, and $\pi_1,\pi_2$ denote the total first and second projections
of the self-delimiting pair encoding, with empty string returned when decoding fails.
\end{theorem}

\begin{theorem}[Quotient and remainder update]
\lean{Complexity.BinArith.div_value_step}
\label{Complexity.BinArith.div_value_step}
\leanok
\difficulty{4}
Let $Q,\rho,v,c,P,T,a\in\mathbb N$ satisfy $c\le1$, $(Qv+\rho)(2P)+(T+cP)=a$, and
$\rho<v$. If $v\le c+2\rho$, then $((1+2Q)v+(c+2\rho-v))P+T=a$ and $c+2\rho-v<v$. If
$v\nleq c+2\rho$, then $(2Qv+(c+2\rho))P+T=a$ and $c+2\rho<v$.
\end{theorem}

\begin{theorem}[Division invariant under one step]
\lean{Complexity.BinArith.divInv_step}
\label{Complexity.BinArith.divInv_step}
\leanok
\uses{Complexity.BinArith.DivInv, Complexity.BinArith.bitsVal_append_singleton, Complexity.BinArith.divStep, Complexity.BinArith.div_value_step, Complexity.BinArith.le, Complexity.BinArith.length_bits_le_of_le, Complexity.BinArith.sub, Complexity.LogProg.bitsVal, Complexity.LogProg.bitsVal_bits, Complexity.isNilB, Complexity.pairFstD, Complexity.pairFstD_pairEncode, Complexity.pairSndD, Complexity.pairSndD_pairEncode}
\difficulty{6}
Let $z$ and $S$ be bit strings, and suppose $S$ satisfies the division invariant for
$z$. Then the state obtained from $S$ by one division step also satisfies that
invariant. If the first projection of the second projection of $S$ is nonempty, the
length of that same nested projection after the step is strictly smaller than before. If
that nested projection is empty, the step leaves $S$ unchanged.
\end{theorem}

\begin{theorem}[Length bound from a division invariant]
\lean{Complexity.BinArith.length_le_of_divInv}
\label{Complexity.BinArith.length_le_of_divInv}
\leanok
\uses{Complexity.BinArith.DivInv, Complexity.length_pair_components_le, Turing.length_pairEncode}
\difficulty{3}
Whenever the binary-division invariant holds for finite Boolean words $z,S$, the state word satisfies $|S|\le7(|z|+1)$.
\end{theorem}

\begin{definition}[Iterated division step]
\lean{Complexity.BinArith.divLoop}
\label{Complexity.BinArith.divLoop}
\leanok
\uses{Complexity.BinArith.divInit, Complexity.BinArith.divStep}
For each finite Boolean sequence $z$, define its division-loop output to be the result
of applying the division step $|z|+1$ times to the division initialization of $z$.
\end{definition}

\begin{theorem}[Polynomial-time computability of the division loop]
\lean{Complexity.BinArith.polyTimeComputable_divLoop}
\label{Complexity.BinArith.polyTimeComputable_divLoop}
\leanok
\uses{Complexity.BinArith.DivInv, Complexity.BinArith.divInit, Complexity.BinArith.divInv_init, Complexity.BinArith.divInv_step, Complexity.BinArith.divLoop, Complexity.BinArith.iterate_invariant, Complexity.BinArith.length_le_of_divInv, Complexity.BinArith.polyTimeComputable_divStep, Complexity.BinArith.polyTimeComputable_iterate, Complexity.BinArith.polyTimeComputable_reverse, Complexity.PolyTimeComputable, Complexity.PolyTimeComputable.comp, Complexity.PolyTimeComputable.pairEncode, Complexity.polyTimeComputable_const, Complexity.polyTimeComputable_pairFstD, Complexity.polyTimeComputable_pairSndD}
\difficulty{3}
The function obtained by applying the division step $|z|+1$ times to the prescribed
division initialization of $z$ is polynomial-time computable.
\end{theorem}

\begin{theorem}[Quotient and remainder of the division loop]
\lean{Complexity.BinArith.divLoop_spec}
\label{Complexity.BinArith.divLoop_spec}
\leanok
\uses{Complexity.BinArith.DivInv, Complexity.BinArith.divInit, Complexity.BinArith.divInv_init, Complexity.BinArith.divInv_step, Complexity.BinArith.divLoop, Complexity.BinArith.divStep, Complexity.BinArith.iterate_invariant, Complexity.BinArith.iterate_measure_eq_zero, Complexity.LogProg.bitsVal, Complexity.length_pairFstD_le, Complexity.pairFstD, Complexity.pairFstD_pairEncode, Complexity.pairSndD, Complexity.pairSndD_pairEncode}
\difficulty{5}
Let $z$ be a bit string whose second projection has positive value. After the division
loop on $z$, the value of the third nested second projection is the quotient of the
value of the first projection by the value of the second projection, and the value of
the first projection of the second nested second projection is the corresponding
remainder.
\end{theorem}

\begin{theorem}[Polynomial-time computability of binary division]
\lean{Complexity.BinArith.polyTimeComputable_div}
\label{Complexity.BinArith.polyTimeComputable_div}
\leanok
\uses{Complexity.BinArith.div, Complexity.BinArith.divLoop_spec, Complexity.BinArith.polyTimeComputable_divLoop, Complexity.BinArith.polyTimeComputable_isZero, Complexity.BinArith.polyTimeComputable_normBits, Complexity.BinArith.polyTimeComputable_slot4, Complexity.LogProg.bitsVal, Complexity.PolyTimeComputable, Complexity.PolyTimeComputable.comp, Complexity.pairSndD, Complexity.polyTimeComputable_const, Complexity.polyTimeComputable_ite, Complexity.polyTimeComputable_pairSndD}
\difficulty{4}
The function on bit strings that returns the binary expansion of the quotient, rounded
down, of the natural numbers represented by the first and second components of the
decoded input is polynomial-time computable.
\end{theorem}

\begin{theorem}[Polynomial-Time Computability of Remainder]
\lean{Complexity.BinArith.polyTimeComputable_mod}
\label{Complexity.BinArith.polyTimeComputable_mod}
\leanok
\uses{Complexity.BinArith.divLoop_spec, Complexity.BinArith.mod, Complexity.BinArith.polyTimeComputable_divLoop, Complexity.BinArith.polyTimeComputable_isZero, Complexity.BinArith.polyTimeComputable_normBits, Complexity.BinArith.polyTimeComputable_slot3, Complexity.LogProg.bitsVal, Complexity.PolyTimeComputable, Complexity.PolyTimeComputable.comp, Complexity.pairFstD, Complexity.pairSndD, Complexity.polyTimeComputable_ite, Complexity.polyTimeComputable_pairFstD, Complexity.polyTimeComputable_pairSndD}
\difficulty{4}
Encode a pair of bit strings by doubling each bit of the first string, appending the
separator $01$, and then appending the second string. Decode aligned pairs as specified
by this encoding, taking a decoded component to be empty if parsing fails, and interpret
each component as a binary numeral with its first bit least significant. The function
sending the encoded input to the binary expansion of the first decoded value modulo the
second decoded value is polynomial-time computable.
\end{theorem}
